\section{Corrections to prior reports}
\label{app:corrections}

Earlier work in this lineage contains measurement and implementation errors that later cycles found.
The project's convention is that these are stated as corrections rather than silently patched into the
artifacts they correct, and the itemised register is placed here so that it does not compete with the
subject of the paper. \S\ref{sec:background} summarises the consequences that bear on
\sestina{}: earlier evaluation lacked the power to resolve the effects it reported on, one
exploitability interpretation was substantially too strong, an exact-belief implementation was unsound
in multi-agent configurations, and the cost of the frontier's search was mischaracterised by two
orders of magnitude. Those discoveries are why the evaluation framework of \S\ref{sec:evaluation}
exists in the form it does.

\subsection{Corrections to earlier FishBot reports}
\label{app:corrections-corpus}

\begin{description}

\item[Search-cost estimates.] The \vsix{} paper recorded the
  searching configuration at ``three orders of magnitude'' the blueprint, mixing an all-threads
  timing with a single-thread one; phase~1 of this cycle re-timed it on one basis at
  \vsevenPhaseOneCost$\times$. Both figures describe the
  \emph{unrestricted} configuration, which is not the point of the frontier anything was measured at.
  The endgame-restricted operating point was re-measured at $\sim$\vsevenCostRecorded$\times$ during
  phase~3, and that correction was carried into the protocol. This report corrects the correction: measured on the shared panel at a fixed thread count,
  \fcheap{}, the configuration the \vsevenCostRecorded$\times$ describes, costs at least
  \vsevenCostFcheapTight$\times$ the blueprint, and \sestina{} at least
  \vsevenCostFrozenTight$\times$ (\S\ref{sec:results-cost}). Every figure in that table is an
  observed whole-match slowdown, diluted by the opponent's cost, and none carries an interval. Both
  earlier errors were ratios quoted without the basis they were computed on.

\item[Interpretation of adversarial-response results.] The v0.4 findings document reports that a policy fitted
  specifically against v0.4 fails to reach parity with it; the same study's own table records the
  response reaching $\vsevenVfourLbr\%$ with an interval spanning parity. The \vsix{} study corrected
  this; it is repeated here because the eight losing searches of \S\ref{sec:results-adversarial} are
  a number of the same kind.

\item[Declaration-allocation comparison.] The register ranked it second by priority, at
  \vsevenLonePriority{}, and estimated it at roughly two percentage points.
  \S\ref{sec:development-alloc} reports that the two allocators selected the same allocation at every
  one of \vsevenLoneDeclsTotal{} recorded declarations (\vsevenLoneDeclsDistinct{} from distinct
  capture cells), that the exact-joint alternative performed worse on the declaration decision, and
  that most of the residual error class is an information limit under the constraint model rather
  than a defect; the tested mechanism therefore produced no difference, which is an empirical
  finding rather than an identity.

\item[Forced-endgame improvement estimate.]
  Re-derived from this cycle's capture, closing the entire accuracy gap is worth
  \vsevenLthirteenClose{}~pp at the measured incidence, the register's figure corrected
  \emph{upward} by a factor of two, and still a fiftieth of the reference bar. The estimate is small
  either way; \S\ref{sec:development-arithmetic} records that the route by which it could grow was
  not tested.

\item[Local-best-response citation correction.] The vetted entry for
  the local-best-response result annotated it as showing ``every ACPC 2016 entrant'' exploitable
  above the stated bound. The source says \emph{each evaluated bot}, and states that the
  competition's winning submission was not evaluated; it also describes the search as looking one
  action ahead, where the annotation said two. Both are corrected in
  \S\ref{sec:related-search} and in the bibliography entry, whose header records the edit.

\item[Detection calibration across evaluation budgets.] Phase~1 measured
  \vsevenFloorPhaseOne{} and predicted the floor would fall as the inverse square root of evaluation
  games. At exactly four times the games the smallest recovered rung was \vsevenPrimaryFloor{}, and
  the \vsevenFloorRungBelow{}~pp rung remained unrecovered, with a phase-2 excess centred near zero
  rather than merely wide. For that responder class and budget the exact result is that two-rung
  outcome; it does not establish that smaller effects are absent. \S\ref{sec:results-controls}
  reports the corresponding controls on holdout material.

\end{description}

\subsection{Corrections to development and evaluation records}
\label{app:corrections-self}


\begin{description}
\item[Repository-history rewrite.] A per-decision
  search capture from the development phase, \vsevenKfiveBytes{} bytes, exceeded the repository
  host's file-size limit and was removed from git history so the work could be published. The file is
  retained on disk with its SHA-256 recorded in a sidecar
  (\texttt{\vsevenKfiveDigest{}\ldots})\@ and is regenerable from the capture harness; no measurement in
  this paper depends on it. The rewrite changed every commit hash from phase 3 onward, so hashes
  recorded inside the phase-5 artifacts no longer resolve.
  Appendix~\ref{app:repro} gives the mapping and confirms the seal commit is unaffected.
\end{description}
\begin{description}

\item[Freeze-verification baseline.] The freeze checker reads its source-drift baseline from the
  freeze artifact's own provenance field, and that artifact was rewritten twice after the freeze
  commit, each rewrite refreshing the baseline to the then-current tree, so the check compared the
  tree against itself and could not fire. Measured against the actual freeze commit,
  \vsevenDriftChanged{} of the \vsevenDriftTotal{} engine sources differ, \vsevenDriftList{}, of
  which \vsevenDriftMidBattery{} arrived in the commits that landed while the battery ran and
  \vsevenDriftPreBattery{} already differed before it began. The policy identity is not weakened:
  the round-trip assertions and the mirror digest are computed by running the engine, and the
  engine rebuilt from the tree as it stood at the end of the battery reproduces the mirror digest
  \texttt{\vsevenMirrorDigest{}}, identical to the frozen artifact and to the binary that played
  every cell (artifact \filepath{P5-drift.json}, deviation D27).

\item[Scored-run reproducibility.]
  \label{app:corrections-reproducibility}
  Before any holdout material was spent, one training cell was checked to be bit-identical over
  \vsevenReproRuns{} runs across three thread counts, and it was
  (\vsevenReproIdentical{}). That check was too small. The panel contains three policies
  twice, so twenty-four pairs of byte-identical commands were run twice each; one pair disagreed, by
  one game in \vsevenResidueGames{}. That is the cross-deal agent residue of \S\ref{sec:results-controls}
  reaching the \emph{scored} mode, which four hundred deals was not large enough to expose. The
  magnitude is roughly \vsevenResidueRatio$\times$ below a cell half-width and no verdict in this paper turns on
  it, but the earlier claim should be read as ``reproducible at four hundred deals'', not
  ``reproducible''.

\item[Experiment-count correction.] The phase-5 recording stated
  \vsevenScoredCellsRecorded{} scored match cells in three places; the figure was hand-typed in the
  reduction script rather than computed. Counting rows that carry a match object across every scored
  battery gives \vsevenScoredCells{}, and the count was corrected from
  \vsevenScoredCellsRecorded{} to \vsevenScoredCells{}. Nothing turned on it, it was the denominator
  of an observation about a single action-limit hit. The generator for this paper computes it, and
  the phase-5 recording has since been corrected to agree.

\item[Partner-comparison uncertainty.]
  \label{app:corrections-intervals}
  The phase-5 results document printed the partner deltas as bare point estimates with a replication
  flag. Every underlying cell carries a deal-clustered interval, and the delta arithmetic is fixed by
  the protocol, so the intervals were computable and had simply been compressed out of a markdown
  table. Table~\ref{tab:partners} carries them, and the phase-5 recording has since been corrected to
  carry them too, at which point the row that decides S1 turned out to have an interval crossing
  its own threshold (\S\ref{sec:results-partners}). The recorded verdict is unchanged; the added
  intervals expose uncertainty not visible in the original point-estimate table.

\item[Cross-play sign convention.] The phase-5 document reports the gap as \vsevenXpGapRecorded{}, which is diagonal minus
  off-diagonal. The claim it settles asks whether the \emph{off}-diagonal collapses, so the informative direction is off-diagonal minus diagonal, which is \vsevenXpGap{}.
  \emph{Neither sign is resolved}: the magnitude is \vsevenXpGapAbs{}~pp against a per-cell
  half-width of \vsevenXpHalfWidth{} and a \vsevenStwoThreshold{}-point threshold, so no verdict
  moves and this paper reads no lift out of it either. What the correction is worth is only this:
  the sign as printed \emph{implies} a small collapse, and the data support no direction at all.

\item[Bibliographic metadata corrections.] During the preparation of the v1.0 report the author
  checked each of the seventy-one references against its source rather than against the earlier
  cycle's bibliography, from which all of them were inherited; this is an internal correction
  inventory, not an independent bibliographic verification, and eight entries carried a metadata
  error. Three author lists
  were wrong: one entry carried the adjacent entry's five names where the paper it cites has four,
  one printed nine names from a later paper by overlapping authors where the cited preprint has
  seven, and one added an author the cited work does not have. The others were a repository start
  year a year early, a page range, the title of a containing work, a web article's site printed as
  its author, and a posting year. All eight are corrected in this document's reference list and none
  of them touches a claim in the text. This revision additionally moved evaluative annotations out
  of five entries and split two entries that referred to distinct documents; the reference list's
  header records those edits.

\item[Rule-variant residual uncertainty.] The phase-5 document forms the residual's half-width
  over five rule-variant cells with unit coefficients, including the \texttt{-{}-maxasks=360} row,
  which is bit-identical to the default on both banks and so carries no variance. Excluding it, and
  giving the default row the coefficient the two-component residual gives it, the interval is
  $[\vsevenLegacyResidualLo, \vsevenLegacyResidualHi]$ rather than the recorded
  $[\vsevenLegacyResidualRecordedLo, \vsevenLegacyResidualRecordedHi]$; this paper prints the
  narrower. No verdict moves: S5 is decided by the largest excursion. Both reported propagated
  ranges include zero, but neither supports a calibrated inferential conclusion: neither is a valid
  95\% interval, because the cells share banks and deals and the ask-cap identity shows the
  correlation is large.

\item[Calibration-rung correspondence.] The
  preregistration introduces B9.3's sub-floor rung as ``\texttt{hstr=0.05}, the rung phase 2 measured
  at $+0.88$'' and directs the phase-5 excess to be read against phase 2's $-0.01$
  $[-0.45, +0.42]$. Phase 2 measured \texttt{hstr=0.05} at a planted cost of $+0.67$ and an excess of
  \vsevenSubFloorExcess{} $[\vsevenSubFloorExcessLo, \vsevenSubFloorExcessHi]$; $+0.88$ and $-0.01$
  are its \texttt{hstr=0.08} row, and phase~1 pins the same split independently. The phase-5
  recording inherited the pairing. \S\ref{sec:results-controls} quotes the rung-matched figure. The
  verdict does not move --- the control is whether the planted edge is recovered, and it is not on
  either pairing --- but the margin it is compared against does, from seven tenths of a point to a
  little under half a point.

\item[Live-ask filtering statistics.] The
  phase-2 mechanism characterisation captures the \emph{target} arm, so every rate in its table is
  measured on \vsix{}. Its rate column comes from a 4,000-deal mirror cell and its
  \texttt{n} column from the 3,000-deal treatment cell, so the pair
  $0.00994$ over $260{,}987$ decisions is not any one cell. The quantity also does not carry over:
  the gate acts on the argmax of an ask score that \sestina{} extends with a coordinate
  \vsix{} does not have, so the rate at which it displaces that argmax is a different number.
  \S\ref{sec:agent-policy} now quotes a capture of the frozen configuration from its own seats on the
  same diagnostic bank, which binds at \vsevenGateBindRate\% of \vsevenGateBindN{} ask decisions,
  half a point above the incumbent's rate.

\item[Duplicate forced-endgame records.] The second sample that refutes the
  one-live-half-suit universal (\S\ref{app:dialect}) is reported in its source as 118 of 120 forced
  declarations. That source's own preceding section establishes that a mirror run at
  \texttt{-{}-rotations=6} plays identical policies in both team orientations and so records every
  declaration twice, and re-runs this very seed at \texttt{-{}-rotations=3} for
  \vsevenLadderSecondTotal{}. The distinct evidence is one occasion in \vsevenLadderSecondTotal{},
  not two in 120. The rate is unchanged; the sample size is halved. The same doubling applies to the
  earlier run quoted in the same sentence, which is why it is described there as records rather than
  as occasions.

\item[Unsupported adversarial-cost estimate.] The phase-2 document, closing the
  forced-endgame induction route, states that the construction which raises the \emph{target's}
  incidence ``was measured at $-17$ points''. No cell in the results tree produces it: no recorded
  configuration string describes that construction, and the only measurement near $-17$ in the cycle
  is an unrelated belief ablation. What was measured is the other construction, a team that never
  declares voluntarily, which raises its own incidence and not its opponent's, at
  \vsevenLthirteenAdvCostLo{} to \vsevenLthirteenAdvCostHi{}~pp on two independent cells.
  \S\ref{sec:development-arithmetic} now says that, and says that the route is untested rather
  than closed by a price; the phase-2 document still carries the unsupported figure.

\item[Source attribution and leaf-evaluation metrics.] The
  \vsevenPhaseOneCost$\times$ cost of the unrestricted search is phase~1 of this cycle's own timing,
  produced expressly to replace the \vsix{} paper's mixed-basis ``three orders of magnitude''; the
  protocol presented it as an inherited record that this cycle corrected. And the leaf evaluator's \vsevenKoneFold-fold gain in between-candidate
  $R^2$ is an in-sample gain, described in its source as replicating out of sample; out of sample the
  same pair is \vsevenKoneRsqOosBefore{} against \vsevenKoneRsqOosAfter{}, so the direction
  replicates and the multiple does not. Neither error changes a verdict.

\end{description}

\subsection{Cross-version comparability}
\label{app:corrections-comparability}

Two things about this cycle's evidence standard differ from earlier cycles', and a reader comparing
across versions should have them. Every claim here is measured on material sealed before the
configuration existed, under a protocol committed before any of it was played, with a stated
failure condition, which no previous study in this corpus did. And every claim is replicated across
two disjoint banks with the sign rule of \S\ref{sec:evaluation-replication} applied, which is why
\S\ref{sec:results-location} reports a non-replication that a pooled-interval convention would have
reported as a pass.
